\documentclass[11pt]{article}

\usepackage[margin=1in]{geometry}
\usepackage[T1]{fontenc}
\usepackage[utf8]{inputenc}
\usepackage{lmodern}
\usepackage{microtype}
\usepackage{amsmath,amssymb,amsthm,mathtools}
\usepackage{bm}
\usepackage{booktabs}
\usepackage{array}
\usepackage{enumitem}
\usepackage{graphicx}
\usepackage[numbers,sort&compress]{natbib}
\usepackage[colorlinks=true,linkcolor=blue!60!black,citecolor=blue!60!black,urlcolor=blue!60!black]{hyperref}

\newtheorem{definition}{Definition}[section]
\newtheorem{proposition}[definition]{Proposition}
\newtheorem{lemma}[definition]{Lemma}
\newtheorem{corollary}[definition]{Corollary}
\theoremstyle{remark}
\newtheorem{remark}[definition]{Remark}
\newtheorem{example}[definition]{Example}

\newcommand{\R}{\mathbb{R}}
\newcommand{\C}{\mathbb{C}}
\newcommand{\HH}{\mathbb{H}}
\newcommand{\SU}{\mathrm{SU}}
\newcommand{\U}{\mathrm{U}}
\newcommand{\Id}{I}
\newcommand{\ii}{\mathbf{i}}
\newcommand{\jj}{\mathbf{j}}
\newcommand{\kk}{\mathbf{k}}
\newcommand{\norm}[1]{\left\lVert #1 \right\rVert}
\newcommand{\abs}[1]{\left\lvert #1 \right\rvert}
\newcommand{\set}[1]{\left\{ #1 \right\}}
\newcommand{\canon}{\operatorname{Can}}
\newcommand{\normalize}{\operatorname{Norm}}
\newcommand{\sgncanon}{\operatorname{SignCan}}
\newcommand{\uoneq}{\mathrm{u1q}}
\newcommand{\firstnz}{\operatorname{firstnz}}

\title{\textbf{A Canonical Quaternionic Intermediate Representation}\\
for Single-Qubit Quantum Circuits}

\author{John G. Van Geem\textsuperscript{*}\\
RQM Technologies\\
orcid: 0009-0002-4003-8452}

\date{Date 2026-04-09}

\begin{document}
\maketitle

\begin{abstract}
This paper defines a canonical quaternionic intermediate representation for
single-qubit circuit segments. Building on the fixed quaternion--$\SU(2)$
correspondence established in \texttt{qqc-001-foundations}, it treats the
physically relevant action of a single-qubit gate modulo global phase by first
choosing an $\SU(2)$ representative and then representing that representative
by a unit quaternion
\[
q = w + x\ii + y\jj + z\kk.
\]
The paper introduces a compiler-facing IR object
\[
\uoneq(w,x,y,z),
\]
formalizes its unit-norm and sign-ambiguity constraints, and defines
deterministic normalization and canonicalization rules. The canonical
representative is chosen from $\set{q,-q}$ by a closed-hemisphere convention
$w>0$, together with an explicit equatorial tie-break when $w=0$. This yields a
shortest-geodesic, deterministic internal form suitable for fusion and rewrite
pipelines. Explicit translations are given for the identity, Pauli gates,
$H$, $S$, $T$ and their adjoints, and axis rotations
$R_x(\theta), R_y(\theta), R_z(\theta)$. Composition of adjacent single-qubit
gates is shown to correspond to quaternion multiplication followed by
re-canonicalization, and worked examples illustrate the relation between named
gate syntax, canonical IR, and single-qubit segments extracted from larger
circuits. The scope is limited to single-qubit gate action modulo global phase;
no multi-qubit IR, optimization guarantee, or hardware claim is made.
\end{abstract}

\noindent\textbf{Keywords:} quaternionic intermediate representation; single-qubit circuits; compiler IR; $\SU(2)$; unit quaternions; circuit canonicalization

\section{Introduction}

The first paper in this series, \texttt{qqc-001-foundations}, established a
conservative mathematical point: after removal of overall phase, a single-qubit
gate admits a representative in $\SU(2)$, and under a fixed standard
convention that representative is modeled exactly by a unit quaternion. The
present paper moves that statement from mathematical language into software
architecture language. Its purpose is to specify a deterministic internal
representation for \emph{single-qubit circuit structure}.

Compiler pipelines need internal forms that are more rigid than frontend
surface syntax. Real circuit languages expose named gates, parameters, control
flow, and scheduling constructs \citep{cross2022,smith2016}. That is
appropriate at the program boundary, but it does not by itself supply a
canonical representation of single-qubit action. Distinct textual gate lists may
denote the same physical action modulo global phase. Likewise, two distinct
$\SU(2)$ representatives may differ only by sign. A simplifier or merger that
operates only on gate names therefore has to recover equivalence information
that is not explicit in the syntax.

The claim of this paper is modest. It does not propose a new quantum language,
and it does not deny the standard matrix formulation of single-qubit gates.
Instead, it defines a canonical intermediate representation layer for
single-qubit gate segments. The representation is quaternionic because the
quaternion model makes composition and sign ambiguity explicit while remaining
exactly equivalent to the standard $\SU(2)$ representative after global phase
has been removed.

The resulting object is called $\uoneq$. It is a quadruple
\[
(w,x,y,z)\in\R^4,\qquad w^2+x^2+y^2+z^2=1,
\]
interpreted as the unit quaternion $q=w+x\ii+y\jj+z\kk$. Canonicalization is
then a separate deterministic step: normalize, resolve the sign ambiguity
$q\sim -q$, and choose a unique representative on a closed hemisphere of $S^3$.
This produces an internal form that is suitable for gate fusion, equality
checks, and segment-level rewrites.

The scope is limited to single-qubit circuits only. Entangling gates appear
only as boundaries that terminate single-qubit segments. No claim is made about
multi-qubit quaternionic ontologies, backend execution, or performance
advantage.

\section{Background and dependency on paper 1}

\subsection{Gates, phase, and the fixed convention}

We carry forward without modification the convention fixed in
\texttt{qqc-001-foundations}. Single-qubit gates live fundamentally in
$\U(2)$. Physical pure-state action is unchanged if a gate is multiplied by an
overall phase. The IR defined here therefore models single-qubit gate action
\emph{modulo global phase}.

\begin{proposition}[Single-qubit gates admit $\SU(2)$ representatives]
For every $U\in\U(2)$ there exist $\chi\in\R$ and $V\in\SU(2)$ such that
\[
U=e^{i\chi}V.
\]
The representative $V$ is unique up to multiplication by $-1$.
\end{proposition}

\begin{proof}
This is the same determinant-removal argument used in
\texttt{qqc-001-foundations}. Choose $\chi$ such that
$e^{2i\chi}=\det U$. Then $V=e^{-i\chi}U$ has determinant $1$, so $V\in\SU(2)$.
Any second determinant-one representative differs from $V$ by a phase whose
square is $1$, hence by $\pm 1$.
\end{proof}

The quaternion convention is also inherited from paper 1:
\begin{equation}
\Phi(a+b\ii+c\jj+d\kk)=
\begin{pmatrix}
a-di & -c-bi \\
c-bi & a+di
\end{pmatrix}
=
a\Id-i(b\sigma_x+c\sigma_y+d\sigma_z).
\label{eq:phi}
\end{equation}
Under this convention the quaternion basis units align directly with the Pauli
axes:
\[
\Phi(\ii)=-i\sigma_x,\qquad \Phi(\jj)=-i\sigma_y,\qquad \Phi(\kk)=-i\sigma_z.
\]

\begin{proposition}[Unit quaternions model $\SU(2)$]
The restriction of $\Phi$ to the unit quaternions
\[
\mathrm{Sp}(1)=\set{q\in\HH:\abs{q}=1}
\]
is a Lie-group isomorphism from $\mathrm{Sp}(1)$ onto $\SU(2)$.
\end{proposition}

\begin{proof}
This is the standard result established in paper 1. The key points are that
$\Phi$ preserves the quaternion multiplication rules and that
$\det \Phi(q)=\abs{q}^2$.
\end{proof}

\begin{corollary}[The only ambiguity is $q\sim -q$]
If $q,q'\in\mathrm{Sp}(1)$ represent the same single-qubit gate action modulo
global phase under the convention \eqref{eq:phi}, then $q'=q$ or $q'=-q$.
\end{corollary}

\begin{proof}
By the previous proposition, unit quaternions correspond bijectively to
$\SU(2)$ elements. By the first proposition, the $\SU(2)$ representative of a
single-qubit gate action is unique up to sign.
\end{proof}

\subsection{Why this matters for compilers}

From a compiler perspective, the point of the previous subsection is simple.
Named gate syntax and canonical IR are different layers.

\begin{itemize}[leftmargin=1.5em]
\item A frontend may say \texttt{H ; T ; Rz(pi/3)} or some equivalent language
surface.
\item The compiler may first translate each gate into a unit quaternion.
\item The compiler may then compose adjacent single-qubit gates by quaternion
multiplication.
\item The compiler may finally store only the resulting canonical
$\uoneq(w,x,y,z)$ tuple.
\end{itemize}

This is not a change of mathematics. It is a change of internal
representation.

\section{Definition of the canonical quaternionic IR}

\begin{definition}[Global-phase equivalence on $\U(2)$]
For single-qubit gates $U,V\in\U(2)$, write
\[
U\sim_{\mathrm{gp}} V
\]
if there exists $\chi\in\R$ such that $U=e^{i\chi}V$.
\end{definition}

\begin{definition}[Quaternionic single-qubit IR element]
A \emph{quaternionic single-qubit IR element}, written
\[
\uoneq(w,x,y,z),
\]
is a quadruple of real numbers satisfying
\[
w^2+x^2+y^2+z^2=1.
\]
It is interpreted as the unit quaternion
\[
q=w+x\ii+y\jj+z\kk\in\mathrm{Sp}(1),
\]
and therefore, under \eqref{eq:phi}, as an $\SU(2)$ representative of a
single-qubit gate action modulo global phase.
\end{definition}

\begin{remark}
The IR stores a representative of the physical single-qubit action, not the
full $\U(2)$ gate with its arbitrary global phase. This is why the natural
ambiguity is the sign pair $\set{q,-q}$ rather than an arbitrary complex
phase.
\end{remark}

\begin{definition}[Equivalence relation on unit quaternions]
For unit quaternions $q,q'\in\mathrm{Sp}(1)$, write
\[
q\sim_{\pm} q'
\]
if $q'=q$ or $q'=-q$.
\end{definition}

\begin{definition}[Single-qubit segment]
A \emph{single-qubit segment} on wire $a$ is a maximal consecutive run of gates
that act only on wire $a$, bounded on either side by the beginning or end of
the circuit, or by an operation that is outside the single-qubit unitary model
for that wire. In particular, an entangling gate touching wire $a$ terminates
the segment.
\end{definition}

\section{Canonicalization rules}

\subsection{Normalization}

Exact mathematics lives on the unit sphere $\mathrm{Sp}(1)\subset\R^4$, but
software implementations often carry floating-point drift.

\begin{definition}[Normalization map]
For any nonzero quadruple
\[
\tilde q=(\tilde w,\tilde x,\tilde y,\tilde z)\in\R^4\setminus\set{0},
\]
define
\[
\normalize(\tilde q)=\frac{\tilde q}{\norm{\tilde q}},
\qquad
\norm{\tilde q}=\sqrt{\tilde w^2+\tilde x^2+\tilde y^2+\tilde z^2}.
\]
\end{definition}

\begin{remark}
The zero quadruple does not represent any gate action and is therefore outside
the domain of canonicalization.
\end{remark}

\subsection{Sign canonicalization}

The quotient ambiguity $q\sim_{\pm}-q$ must be resolved deterministically.

\begin{definition}[Canonical hemisphere]
Let $\mathcal{H}_{\mathrm{can}}\subset \mathrm{Sp}(1)$ be the set of unit
quaternions satisfying either
\begin{enumerate}[label=(\alph*), leftmargin=1.5em]
\item $w>0$, or
\item $w=0$ and the first nonzero component among $(x,y,z)$ is positive.
\end{enumerate}
\end{definition}

\begin{definition}[Sign-canonicalization map]
For $q=(w,x,y,z)\in\mathrm{Sp}(1)$, define $\sgncanon(q)$ by
\[
\sgncanon(q)=
\begin{cases}
q, & w>0,\\
-q, & w<0,\\
q, & w=0 \text{ and the first nonzero entry of } (x,y,z) \text{ is positive},\\
-q, & w=0 \text{ and the first nonzero entry of } (x,y,z) \text{ is negative}.
\end{cases}
\]
\end{definition}

\begin{definition}[Canonical representative]
For any nonzero quaternionic quadruple $\tilde q\in\R^4\setminus\set{0}$,
define its canonical representative by
\[
\canon(\tilde q)=\sgncanon(\normalize(\tilde q)).
\]
If $q\in\mathrm{Sp}(1)$ is already unit norm, then $\canon(q)=\sgncanon(q)$.
\end{definition}

\begin{proposition}[Canonicalization selects one representative from each sign class]
Every equivalence class $\set{q,-q}$ with $q\in\mathrm{Sp}(1)$ contains exactly
one element of $\mathcal{H}_{\mathrm{can}}$.
\end{proposition}

\begin{proof}
If $w\neq 0$, then exactly one of $q$ or $-q$ has positive real part. If
$w=0$, then $q$ and $-q$ have opposite signs in the first nonzero component of
$(x,y,z)$, so exactly one satisfies the equatorial tie-break rule. Hence each
class intersects $\mathcal{H}_{\mathrm{can}}$ in exactly one point.
\end{proof}

\begin{remark}[Shortest-geodesic convention]
Write a unit quaternion in axis--angle form as
\[
q=\cos(\theta/2)+\sin(\theta/2)(n_x\ii+n_y\jj+n_z\kk).
\]
The condition $w=\cos(\theta/2)\geq 0$ implies $\theta\in[0,\pi]$. The
canonical hemisphere therefore chooses the representative whose rotation angle
lies on the shorter closed interval. When $w=0$, corresponding to $\theta=\pi$,
both signs have the same geodesic length, so the equatorial tie-break supplies
determinism rather than new physics.
\end{remark}

\begin{remark}[Exact versus numerical equality]
The mathematical equivalence relation is exact. In floating-point
implementations, comparisons should be performed with explicit tolerances after
normalization and canonicalization. That is an implementation detail, not part
of the exact equivalence relation itself.
\end{remark}

\section{Translation from named gates into quaternion IR}

Named gate syntax is not itself canonical. Different frontends may expose the
same physical action with different symbols, parameters, or decompositions. The
translation layer maps those named gates into a common quaternionic form.

\begin{proposition}[Standard named gates have explicit canonical $\uoneq$ entries]
Under the convention \eqref{eq:phi}, the following single-qubit gates map to
the indicated canonical unit quaternions modulo global phase:
\begin{center}
\begin{tabular}{@{}lll@{}}
\toprule
Gate & Chosen representative modulo global phase & Canonical quaternion \\
\midrule
$\Id$ & $\Id$ & $(1,0,0,0)$ \\
$X$ & $-iX$ & $(0,1,0,0)$ \\
$Y$ & $-iY$ & $(0,0,1,0)$ \\
$Z$ & $-iZ$ & $(0,0,0,1)$ \\
$H$ & $-iH$ & $(0,\tfrac{1}{\sqrt2},0,\tfrac{1}{\sqrt2})$ \\
$S$ & determinant-one representative of $S$ & $(\tfrac{1}{\sqrt2},0,0,\tfrac{1}{\sqrt2})$ \\
$S^\dagger$ & determinant-one representative of $S^\dagger$ & $(\tfrac{1}{\sqrt2},0,0,-\tfrac{1}{\sqrt2})$ \\
$T$ & determinant-one representative of $T$ & $(\cos\tfrac{\pi}{8},0,0,\sin\tfrac{\pi}{8})$ \\
$T^\dagger$ & determinant-one representative of $T^\dagger$ & $(\cos\tfrac{\pi}{8},0,0,-\sin\tfrac{\pi}{8})$ \\
$R_x(\theta)$ & $R_x(\theta)$ & $(\cos\tfrac{\theta}{2},\sin\tfrac{\theta}{2},0,0)$ \\
$R_y(\theta)$ & $R_y(\theta)$ & $(\cos\tfrac{\theta}{2},0,\sin\tfrac{\theta}{2},0)$ \\
$R_z(\theta)$ & $R_z(\theta)$ & $(\cos\tfrac{\theta}{2},0,0,\sin\tfrac{\theta}{2})$ \\
\bottomrule
\end{tabular}
\end{center}
\end{proposition}

\begin{proof}
The rotation families follow immediately from the quaternion axis--angle form
\[
q=\cos(\theta/2)+\sin(\theta/2)(n_x\ii+n_y\jj+n_z\kk).
\]
The Pauli gates are the $\pi$-rotations about the coordinate axes, so their
canonical representatives are $\ii,\jj,\kk$. The $S$ and $T$ families are the
$z$-axis quarter- and eighth-turns after removal of determinant phase. For the
Hadamard gate,
\[
H=\frac{\sigma_x+\sigma_z}{\sqrt2},
\]
so after removing a phase factor one obtains the $\pi$-rotation about the axis
$(1,0,1)/\sqrt2$, giving
$q_H=(0,1/\sqrt2,0,1/\sqrt2)$.
\end{proof}

\begin{remark}
The table above concerns semantic gate action modulo global phase. It is
therefore a representation table, not a commitment to any specific frontend
syntax spelling. For example, \texttt{sdg}, $S^\dagger$, or an equivalent
phase-gate spelling may all map to the same canonical tuple.
\end{remark}

\section{Composition and segment representation}

The group law on unit quaternions makes segment fusion immediate.

\begin{proposition}[Composition of adjacent single-qubit gates]
Let $g_1,\dots,g_m$ be consecutive single-qubit gates acting on the same wire,
listed in application order. Let $q_1,\dots,q_m$ be their unit-quaternion
representatives under \eqref{eq:phi}. Then the fused single-qubit action of
the segment is represented by
\[
q_{\mathrm{seg}}=\canon(q_m q_{m-1}\cdots q_1).
\]
\end{proposition}

\begin{proof}
By the convention carried forward from paper 1,
\[
\Phi(q_m q_{m-1}\cdots q_1)
=
\Phi(q_m)\Phi(q_{m-1})\cdots\Phi(q_1),
\]
which is exactly the $\SU(2)$ representative of sequential application of the
segment gates. Canonicalization then chooses the deterministic representative of
the resulting sign class.
\end{proof}

\begin{corollary}[Canonicalization may be applied after each fusion step]
In exact arithmetic, repeated canonicalization during a left-to-right or
right-to-left fold produces the same physical single-qubit action as
canonicalizing once at the end. In floating-point arithmetic, intermediate
re-normalization may improve stability without changing the intended exact
semantics.
\end{corollary}

\begin{remark}
This is the principal compiler-facing advantage of the representation. A
frontend may emit a list of named gates, but an optimizer may replace the list
by a single canonical $\uoneq$ payload whenever a maximal single-qubit segment
is identified.
\end{remark}

\section{Worked examples}

The examples below use a semicolon-separated gate-list style close to common
quantum assembly notations \citep{cross2022,smith2016}. The point is not to
fix a syntax standard, but to show the translation boundary between named gates
and canonical IR.

\begin{example}[A mixed named and parametric segment]
Consider the single-qubit segment
\[
H\;;\;T\;;\;R_z(\pi/3).
\]
The gatewise quaternion entries are
\[
q_H=\left(0,\frac{1}{\sqrt2},0,\frac{1}{\sqrt2}\right),\qquad
q_T=\left(\cos\frac{\pi}{8},0,0,\sin\frac{\pi}{8}\right),
\]
\[
q_{R_z(\pi/3)}=\left(\frac{\sqrt3}{2},0,0,\frac12\right).
\]
Composing in application order gives
\[
q_{\mathrm{seg}}=
\canon\!\bigl(q_{R_z(\pi/3)} q_T q_H\bigr)
\approx
(0.5609855,\,-0.4304593,\,-0.5609855,\,-0.4304593).
\]
The named syntax and the canonical tuple are distinct objects: the former is a
frontend presentation, while the latter is a deterministic segment payload.
\end{example}

\begin{example}[Fusion of two axis rotations]
For
\[
R_x(\pi/2)\;;\;R_y(\pi/2),
\]
the quaternions are
\[
q_x=\left(\frac{\sqrt2}{2},\frac{\sqrt2}{2},0,0\right),\qquad
q_y=\left(\frac{\sqrt2}{2},0,\frac{\sqrt2}{2},0\right).
\]
Hence
\[
\canon(q_y q_x)=\left(\frac12,\frac12,\frac12,-\frac12\right).
\]
This exact tuple is often easier to compare and store than the original
rotation syntax.
\end{example}

\begin{example}[Named syntax can collapse to a simpler canonical payload]
Consider
\[
Z\;;\;S\;;\;T.
\]
The gatewise representatives are
\[
q_Z=(0,0,0,1),\qquad
q_S=\left(\frac{1}{\sqrt2},0,0,\frac{1}{\sqrt2}\right),\qquad
q_T=\left(\cos\frac{\pi}{8},0,0,\sin\frac{\pi}{8}\right).
\]
The fused canonical result is
\[
\canon(q_T q_S q_Z)=\left(\cos\frac{\pi}{8},0,0,-\sin\frac{\pi}{8}\right),
\]
which is the canonical representative of the same single-qubit action modulo
global phase. The example illustrates why a canonical IR is more informative
for simplification than the original named gate list.
\end{example}

\begin{example}[Segment extraction inside a larger circuit]
Consider the two-wire circuit
\[
\text{wire 0: } H\;;\;\mathrm{CNOT}_{0,1}\;;\;T^\dagger\;;\;R_x(\pi/4),
\]
\[
\text{wire 1: } S\;;\;\mathrm{CNOT}_{0,1}\;;\;H.
\]
The entangling gate is outside the present IR. Its only role here is to split
single-qubit segments.

On wire $0$, the segment before the entangling boundary is the single gate
$H$, while the segment after the boundary is
\[
T^\dagger\;;\;R_x(\pi/4),
\]
with fused canonical payload
\[
\canon\!\bigl(q_{R_x(\pi/4)} q_{T^\dagger}\bigr)
\approx
(0.8535534,\,0.3535534,\,0.1464466,\,-0.3535534).
\]
On wire $1$, the entangling gate similarly separates the one-gate segments
$S$ and $H$. The example is intentionally modest: the paper does not define the
internal representation of the entangling gate, only the single-qubit segments
around it.
\end{example}

\section{Standard mathematics vs.\ RQM Technologies framing}

\subsection{Standard facts used here}

The following ingredients are standard and not specific to this series:
\begin{enumerate}[label=(\arabic*), leftmargin=1.5em]
\item single-qubit gates live in $\U(2)$;
\item after removing global phase, they admit representatives in $\SU(2)$;
\item the unit quaternions form a Lie group isomorphic to $\SU(2)$;
\item the only residual ambiguity in that model is $q\sim -q$; and
\item sequential single-qubit gate composition corresponds to quaternion
multiplication under the fixed convention \eqref{eq:phi}.
\end{enumerate}

\subsection{Project-specific design choices}

The project-specific content of this paper is representational rather than
physical:
\begin{enumerate}[label=(\alph*), leftmargin=1.5em]
\item the IR object name $\uoneq$;
\item the decision to store single-qubit action as a four-real quaternion
tuple;
\item the closed-hemisphere sign convention $w>0$ plus the equatorial tie-break
rule; and
\item the interpretation of maximal single-qubit runs between entangling
boundaries as the natural unit of fusion for this layer.
\end{enumerate}

These are software-architecture choices. They are intended to make equality,
fusion, and simplification deterministic. They are not claims about new
physics, and they do not imply that named gate syntax or matrix methods are
incorrect.

\section{Scope, limitations, and non-claims}

This paper makes several non-claims explicit:
\begin{enumerate}[label=(\arabic*), leftmargin=1.5em]
\item It does \emph{not} define a full multi-qubit quaternionic circuit IR.
\item It does \emph{not} specify the internal semantics of entangling gates.
\item It does \emph{not} prove a runtime, memory, or optimization advantage.
\item It does \emph{not} replace existing SDKs, assembly languages, or compiler
frameworks.
\item It does \emph{not} define a backend execution format.
\item It does \emph{not} claim that standard matrix methods are wrong or
insufficient.
\end{enumerate}

The paper defines only a canonical representation layer for single-qubit gate
action modulo global phase.

\section{Conclusion}

The main result of this paper is representational. A single-qubit gate action
modulo global phase can be stored deterministically as a canonical unit
quaternion
\[
\uoneq(w,x,y,z),
\]
with explicit normalization and sign-resolution rules. The representation is
equivalent to the standard $\SU(2)$ representative after global phase removal,
inherits a simple composition law from quaternion multiplication, and provides a
clean internal target for single-qubit segment fusion.

In that sense, the paper is a bridge from quaternionic single-qubit geometry to
compiler IR design. It does not solve all compiler problems, but it does make
one layer precise: the canonical internal representation of single-qubit gate
segments.

\appendix

\section{Lookup table for common gates}

\begin{center}
\begin{tabular}{@{}lll@{}}
\toprule
Gate & Canonical quaternion & Notes \\
\midrule
$\Id$ & $(1,0,0,0)$ & identity \\
$X$ & $(0,1,0,0)$ & $\pi$-rotation about $x$ \\
$Y$ & $(0,0,1,0)$ & $\pi$-rotation about $y$ \\
$Z$ & $(0,0,0,1)$ & $\pi$-rotation about $z$ \\
$H$ & $(0,\tfrac{1}{\sqrt2},0,\tfrac{1}{\sqrt2})$ & $\pi$-rotation about $(1,0,1)/\sqrt2$ \\
$S$ & $(\tfrac{1}{\sqrt2},0,0,\tfrac{1}{\sqrt2})$ & quarter-turn about $z$ \\
$T$ & $(\cos\tfrac{\pi}{8},0,0,\sin\tfrac{\pi}{8})$ & eighth-turn about $z$ \\
$R_x(\theta)$ & $(\cos\tfrac{\theta}{2},\sin\tfrac{\theta}{2},0,0)$ & exact family \\
$R_y(\theta)$ & $(\cos\tfrac{\theta}{2},0,\sin\tfrac{\theta}{2},0)$ & exact family \\
$R_z(\theta)$ & $(\cos\tfrac{\theta}{2},0,0,\sin\tfrac{\theta}{2})$ & exact family \\
\bottomrule
\end{tabular}
\end{center}

\section{Implementation notes}

The mathematical definition of $\canon$ is exact, but an implementation should
usually include three practical checks:
\begin{enumerate}[label=(\alph*), leftmargin=1.5em]
\item reject or special-case near-zero input before normalization;
\item re-normalize after quaternion multiplication to control floating-point
drift; and
\item evaluate the $w=0$ tie-break with a tolerance so that numerically small
values are not misclassified.
\end{enumerate}

These are engineering details layered on top of the exact canonicalization rule.

\section{Companion figures}

Publication-ready SVG figures for the package are provided in
\texttt{artifacts/figures/}. They illustrate the frontend-to-IR translation
pipeline, the quotient identification $q\sim -q$, the canonical hemisphere
selection rule, and segment extraction around entangling boundaries.

\bibliographystyle{plainnat}
\bibliography{references}

\end{document}
